\documentclass[10pt,a4paper]{amsart}
\usepackage[margin=19mm]{geometry}
\usepackage{amsmath,amssymb,mathtools}
\usepackage[scr=rsfs,cal=euler]{mathalfa}
\usepackage{xfrac}
\usepackage[unicode,hidelinks,bookmarks=false]{hyperref}
\newcommand{\ie}{i.e.\ }
\newcommand{\cf}{cf.\ }
\newcommand{\resp}{respectively\ }
\newcommand{\Subsection}{\subsection}
\providecommand{\nobreakdash}{\nobreak\hbox{-}}
\setlength{\emergencystretch}{2em}
\multlinegap0pt
\usepackage{amssymb}
\usepackage[shortlabels]{enumitem}
\usepackage{mathtools}
\usepackage{tikz}
\newtheorem{lemma}{Lemma}
\title{Existence, uniqueness and regularity of solutions to the parabolic Ambrosio-Tortorelli system}
\author{Martin Rakovsky}
\date{Source: arXiv:2607.03373v1 (CC BY 4.0)}
\begin{document}
\maketitle

\noindent\emph{Source and licence.} Existence, uniqueness and regularity of solutions to the parabolic Ambrosio-Tortorelli system. Version arXiv:2607.03373v1 (CC BY 4.0), distributed by arXiv under the Creative Commons Attribution 4.0 International licence, http://creativecommons.org/licenses/by/4.0/, as recorded in the arXiv licence metadata for this exact version and shown on the abstract page https://arxiv.org/abs/2607.03373v1. Full licence notice: \texttt{licenses/arxiv-2607.03373-CC-BY-4.0.txt}.\par\smallskip

\noindent\textit{Editorial scope.} Counted unit: the article's lemma equ (source0.tex lines 1737--1753). The article's complete original proof of it is delivered with it (source0.tex lines 1755--1817, one \texttt{proof} environment of 63 lines). No mathematics is written, repaired or re-derived: the statement and the proof are the article's own lines, in the article's own notation, and no retained line is substituted or re-flowed.  Retained uncounted as the source-local context the proof needs: lines 1565--1572; lines 1576--1578; lines 1711--1716; lines 1722--1727.  Nothing was omitted from any retained span, so the package carries no omission record.  No retained line contains a citation, so the wrapper carries no bibliography and no bibliography entry.  No retained line contains a figure environment, an image-inclusion command or drawing code, so no image is delivered and the package declares no asset dependency.  Editorial formatting only, in addition: an AMS \texttt{amsart} preamble that restates the article's own declarations and macros used by the retained lines, namely the environments \texttt{lemma} on separate counters, together with the standard packages those lines need.  The retained environments print in this document as Lemma 1 in order of appearance, whereas the article prints the same items with the numbering of its own full document.  The complete native source of the article as distributed in the arXiv e-print for this exact version is bundled unchanged as \texttt{sources/204378/source0.tex}.
\begin{equation}\label{extension uv eq}
\left\{\begin{array}{ll}
\partial_t \overline{u} - \text{div}((\eta_\varepsilon + \overline{v}^2) \nabla \overline{u}) = - \text{div}((\eta_\varepsilon + v_0^2) \nabla u_0) 1_{\{t\leqslant 1\}}, &\text{in } \mathcal{D}'((0,\infty)\times\Omega)\\
\partial_t \overline{v} - \varepsilon \Delta \overline{v} + \frac{\overline{v}-1}{4\varepsilon} + \overline{v}|\nabla \overline{u}|^2 = (- \varepsilon \Delta v_0 + \frac{v_0-1}{4\varepsilon} + v_0|\nabla u_0|^2)1_{\{t\leqslant 1\}},  &\text{in } \mathcal{D}'((0,\infty)\times \Omega)\\
\overline{v}(t,x) = 1 & \text{on } (0,\infty) \times \partial\Omega \\
\overline{u}(t,x) = g(t-1,x)1_{\{t > 1\}} + g(0,x) 1_{\{t\leqslant 1\}} & \text{on } (0,\infty) \times \partial\Omega\\
\end{array}\right. 
\end{equation}

\begin{equation}\label{def U0 V0}
U_0 = - \text{div}((\eta_\varepsilon + v_0^2) \nabla u_0) 1_{\{t\leqslant 1\}}, \quad V_0 = \left(- \varepsilon \Delta v_0 + \frac{v_0-1}{4\varepsilon} + v_0|\nabla u_0|^2\right)1_{\{t\leqslant 1\}}.
\end{equation}

\begin{equation}\label{def : u chapeau}
(\hat{u}(t,x),\hat{v}(t,x))  := \left\{ \begin{array}{ll}
(\overline{u}(t,x), \overline{v}(t,x))  & \text{ if } x\in \widetilde{\Omega} \cap \Omega \\
(\overline{u} (t,\sigma (x)), \overline{v} (t,\sigma(x)) & \text{ if } x\in \widetilde{\Omega} \setminus \Omega,\\
\end{array} \right.
\end{equation}

\begin{equation}\label{def : u tilde}
(\widetilde{u}(t,x), \widetilde{v}(t,x)) := \left\{ \begin{array}{ll}
(\overline{u} (t,x), \overline{v}(t,x)) & \text{ if } x\in \widetilde{\Omega} \cap \Omega \\
(2 \overline{g}(t,\pi(x))- \overline{u} (t,\sigma (x)), 2-\overline{v}(t,\sigma(x))) & \text{ if } x\in \widetilde{\Omega} \setminus \Omega,\\
\end{array} \right.
\end{equation}

\begin{lemma}
Recall the definition \eqref{def U0 V0} of $(U_0, V_0)$
The functions $\widetilde{u}, \widetilde{v}, \hat{u}, \hat{v}$ as defined in \eqref{def : u chapeau} and \eqref{def : u tilde} satisfy 

\begin{multline}\label{boundary : equ}
j\, \partial_t \widetilde{u} - div((\eta_\varepsilon + \hat{v}^2_\varepsilon ) M\nabla \widetilde{u} )\\
 = 2j\, \partial_t \widetilde{g} 1_{\widetilde{\Omega}\setminus \overline{\Omega}} - 2 div(1_{\widetilde{\Omega}\setminus \overline{\Omega}} (\eta_\varepsilon + \hat{v}^2) M \nabla \widetilde{g})+ \frac{j}{2}(U_0 1_{\widetilde{\Omega}\cap \Omega} - U_0\circ \sigma  1_{\widetilde{\Omega}\setminus \Omega}) \quad \text{in } \mathcal{D}'((0,\infty) \times \widetilde{\Omega}),
\end{multline}
%\text{in } \mathcal{D}'((0,+\infty)\times \widetilde{\Omega})


\begin{multline}\label{boundary : eqv}
j\, \partial_t \widetilde{v} - \text{div}(M\nabla \widetilde{v})\\
 =(1_{\widetilde{\Omega} \cap \Omega} - 1_{\widetilde{\Omega}\setminus \overline{\Omega}}) \left( \frac{j}{4\varepsilon} (1-\hat{v}) - (M \widetilde{v} \nabla \widetilde{u} \cdot \nabla \widetilde{u})\right) + \frac{j}{2}(V_0 1_{\widetilde{\Omega}\cap \Omega} - V_0\circ \sigma  1_{\widetilde{\Omega}\setminus \Omega}) \quad \text{in } \mathcal{D}'((0,\infty) \times \widetilde{\Omega}).
\end{multline}
%\text{in } \mathcal{D}'((0,+\infty)\times \widetilde{\Omega}).
\end{lemma}

\begin{proof}
Let $\varphi \in \mathcal{C}_c^\infty ((0,\infty)\times \widetilde{\Omega})$. We decompose $\varphi = \varphi ^s +\varphi^a$ as the sum of its symmetric and anti-symmetric part, with 

\[\varphi^s(x) = \frac{1}{2}( \varphi(x) + \varphi (\sigma (x))), \quad \text{ and } \quad \varphi^a(x) = \frac{1}{2}( \varphi(x) - \varphi (\sigma (x))).\]
The functions $\varphi^s$ and $\varphi ^a$ belong to $\mathcal{C}^{1,1}((0,\infty)\times \Omega))$ and satisfy $\varphi ^s(t,\sigma(x)) = \varphi ^s(t,x)$ and $\varphi ^a (t, \sigma(x)) = - \varphi ^a(t,x)$ for all $(t,x) \in (0,\infty) \times \widetilde{\Omega}$. 

\medskip

\underline{Proof of \eqref{boundary : equ} :} In the expression   

\[\int_0^\infty \int_{\widetilde{\Omega}} j \partial_t \widetilde{u} \varphi + \int_{\widetilde{\Omega}} (\eta_\varepsilon + \hat{v}^2 )M \nabla \widetilde{u} \cdot \nabla \varphi,\]
we consider separately the terms with $\varphi^s$ and the terms with $\varphi^a$ : 
\begin{multline}\label{boundary equ 1}
\int_0^\infty\int_{\widetilde{\Omega}\setminus \overline{\Omega}} j \partial_t \widetilde{u}\: \varphi^s  + \int_0^\infty \int_{\widetilde{\Omega}\setminus \overline{\Omega}} (\eta_\varepsilon + \hat{v}^2) M\nabla \widetilde{u} \cdot \nabla \varphi^s \\
= 2 \int_0^\infty \int_{\widetilde{\Omega}\setminus \overline{\Omega}} \left[j\partial_t \widetilde{g} \varphi^s +(\eta_\varepsilon + \hat{v}^2)M\nabla \widetilde{g} \cdot \nabla \varphi ^s\right]\\
- \int_{\widetilde{\Omega}\setminus \overline{\Omega}} \left[j\partial_t (\overline{u}\, \circ \, \sigma) \varphi^s + (\eta_\varepsilon + \hat{v}^2)M\nabla (\overline{u}\, \circ\, \sigma) \cdot \nabla \varphi ^s\right].
\end{multline}
Using $\varphi^s \circ \sigma = \varphi^s$ and changing the variable : 

\begin{multline}\label{boundary equ 2}
\int_0^\infty\int_{\widetilde{\Omega}\setminus \overline{\Omega}}\left[ j(\partial_t\overline{u}\, \circ \,\sigma)(\varphi^s \circ\, \sigma) + (\eta_\varepsilon + \hat{v}^2)M\nabla (\overline{u}\, \circ \,\sigma) \cdot \nabla (\varphi ^s \circ \,\sigma)\right]\\
= \int_{\widetilde{\Omega}\cap \Omega} \left[\partial_t \overline{u} \varphi^s + (\eta_\varepsilon + \overline{v}^2) \nabla \overline{u} \cdot \nabla \varphi^s\right].
\end{multline}
From \eqref{boundary equ 1} and \eqref{boundary equ 2}, we deduce 

\begin{equation}\label{boundary equ 3}
\int_0^\infty \int_{\widetilde{\Omega}}\left[ j \partial_t \widetilde{u} \,\varphi^s + (\eta_\varepsilon + \hat{v}^2) M \nabla \widetilde{u} \cdot \nabla \varphi^s\right] =  2\int_0^\infty \int_{\widetilde{\Omega}\setminus \overline{\Omega}}\left[ j\partial_t \widetilde{g} \,\varphi^s +(\eta_\varepsilon + \hat{v}^2)M\nabla \widetilde{g} \cdot \nabla \varphi ^s\right].
\end{equation}
We now focus on the part with $\varphi^a$. Notice that $\varphi^a =0$ on $(0,\infty) \times \partial\Omega$, so that  

\begin{multline}\label{boundary equ 4}
\int_0^\infty\int_{\widetilde{\Omega}\setminus \overline{\Omega}} j \partial_t \widetilde{u}\, \varphi^a  + \int_0^\infty\int_{\widetilde{\Omega}\setminus \overline{\Omega}} (\eta_\varepsilon + \hat{v}^2 ) M\nabla \widetilde{u} \cdot \nabla \varphi^a \\
= 2\int_0^\infty \int_{\widetilde{\Omega}\setminus \overline{\Omega}} \left[ j\partial_t \widetilde{g} \varphi^a +(\eta_\varepsilon + \hat{v}^2)M\nabla \widetilde{g} \cdot \nabla \varphi ^a \right]\\
 - \int_0^\infty\int_{\widetilde{\Omega}\setminus \overline{\Omega}} \left[ j\partial_t (u\, \circ \,\sigma) \varphi^a + (\eta_\varepsilon + \hat{v}^2)M\nabla (u\, \circ \,\sigma) \cdot \nabla \varphi ^a \right].
\end{multline}
Using that $\varphi^a \circ \sigma = -\varphi^a$ and changing the variable : 

\begin{align*}
&\int_0^\infty\int_{\widetilde{\Omega}\setminus \overline{\Omega}} \left[j\partial_t (u\, \circ \,\sigma) \varphi^a + (\eta_\varepsilon + \hat{v}^2)M\nabla (u\, \circ \,\sigma) \cdot \nabla \varphi ^a\right]\\
&= -\int_0^\infty\int_{\widetilde{\Omega}\setminus \overline{\Omega}}\left[ j\partial_t (u\, \circ \,\sigma) (\varphi^a \circ \sigma) + (\eta_\varepsilon + \hat{v}^2)M\nabla (u\, \circ \,\sigma) \cdot \nabla (\varphi ^a \circ \sigma)\right]\\
&= -\int_0^\infty\int_{\widetilde{\Omega}\cap \Omega}\left[ \partial_t \overline{u} \varphi^a + (\eta_\varepsilon + \overline{v}^2) \nabla \overline{u} \cdot \nabla \varphi^a\right] \\
&=-\int_0^\infty\int_{\widetilde{\Omega}\cap \Omega} U_0 \varphi^a,\\
\end{align*}
where we used \eqref{extension uv eq} in the last equality. 
%&=- \int_{\widetilde{\Omega}\setminus \overline{\Omega}}\left[ j(x)\partial_t u (\sigma) \varphi^a (\sigma) + (\eta_\varepsilon + \hat{v}^2)\nabla u (\sigma) \cdot \nabla \varphi ^a (\sigma) \, j(x)\right]\\
We deduce from \eqref{boundary equ 4} that 

\begin{align}\label{boundary equ 5}
&\int_0^\infty\int_{\widetilde{\Omega}} j\, \partial_t \widetilde{u}\, \varphi^a + (\eta_\varepsilon + \hat{v}^2) M \nabla \tilde{u} \cdot \nabla \varphi^a \nonumber\\
&= \int_0^\infty\int_{\widetilde{\Omega}\setminus \overline{\Omega}} j\, \partial_t \widetilde{u} \varphi^a  + \int_0^\infty\int_{\widetilde{\Omega}\setminus \overline{\Omega}} (\eta_\varepsilon + \hat{v}^2_\varepsilon ) M\nabla \widetilde{u} \cdot \nabla \varphi^a \nonumber\\
&\quad +\int_0^\infty\int_{\widetilde{\Omega}\cap \Omega} \left[\partial_t \overline{u} \varphi^a + (\eta_\varepsilon + \overline{v}^2) \nabla u \cdot \nabla \varphi^a \right] \nonumber\\
&= 2\int_0^\infty\int_{\widetilde{\Omega}\setminus \overline{\Omega}}\left[ j\,\partial_t \widetilde{g} \,\varphi^a +(\eta_\varepsilon + \hat{v}^2)M\nabla \widetilde{g} \cdot \nabla \varphi ^a\right] + \int_0^\infty\int_{\widetilde{\Omega} \cap \Omega} U_0 \varphi^a + 0. \nonumber\\
\end{align}
We observe that
\begin{multline}\label{eq U0nouveau}
\int_0^\infty\int_{\widetilde{\Omega}\cap \Omega} U_0 \varphi^a = \frac{1}{2}\int_0^\infty \int_{\widetilde{\Omega}\cap \Omega} U_0(\varphi - \varphi\circ \sigma)\\
= \frac{1}{2}\int_0^\infty\int_{\widetilde{\Omega}\cap \Omega} U_0\varphi - \frac{1}{2}\int_0^\infty\int_{\widetilde{\Omega}\setminus \Omega} U_0\circ \sigma \varphi j = \int_0^\infty \int_{\widetilde{\Omega}} \frac{j}{2}(U_0 1_{\widetilde{\Omega}\cap \Omega} - U_0\circ \sigma  1_{\widetilde{\Omega}\setminus \Omega})\varphi.
\end{multline}
%

Combining \eqref{boundary equ 3},  \eqref{boundary equ 5} and \eqref{eq U0nouveau} , we obtain \eqref{boundary : equ}. 
The same computation holds for \eqref{boundary : eqv}. 
\end{proof}

\end{document}
